\section{Discussion}
\label{sec:discussion}

\subsection{Routing is a disclosure and latency mechanism, not a green one}

The case for local--edge routing in a screening pipeline is sound, but it is a
different case from the one usually made, and narrower than we expected to be
able to make. On our evaluation the mechanism buys two things a laboratory can
verify directly, and both survive antigen-level resampling. Screening finishes
in roughly a third of the time always-remote execution takes. Most antigens are
never transmitted. The second gain is not obtainable by buying a faster remote
service.

What it does not buy, on evidence of this size, is a demonstrated quality
position. Against always-remote execution the quality difference is not
resolved; against always-local execution only one metric of four separates
them. We had originally read the same comparisons as decisive, on intervals
computed over seeds, and Section~\ref{sec:routing-results} explains why those
were the wrong intervals. A laboratory deciding whether to adopt routing should
therefore treat the latency and disclosure gains as established here and the
quality question as open. That is a less satisfying recommendation than the one
we set out to make. It is also a more defensible one.

What the mechanism does not buy is lower energy, and
Equation~\ref{eq:carbon-threshold} says precisely when it would. Offloading
saves energy only above a ratio of work offloaded to overhead incurred
\citep{kumar2010cloud}. Our contribution is to locate that boundary for this
task, at a routed fraction of \CarbonBreakEvenRate{}, and to observe that a
policy tuned for predictive quality routes about twice as often as that. The
green argument
would recover for longer antigens or a heavier remote tier, where the ratio
shifts. Practitioners can evaluate Equation~\ref{eq:carbon-threshold} for their
own pair rather than assuming either direction.

The practical recommendation, then, is to adopt routing when sequence disclosure
or turnaround is the binding constraint, and not to adopt it on environmental
grounds without measuring the ratio for the models actually in use.

\subsection{What the fusion result means, and what it does not}

Our result is that a convex blend of CARBON-BCE-R with BepiPred-3.0 shows no
advantage over BepiPred-3.0 alone, on either of two datasets, once the blend is
chosen without reference to the evaluation labels. The scope of that claim
matters.

It is not a claim that score fusion is useless in epitope prediction. It is a
claim about one pair of predictors on two splits. One member of that pair is
considerably weaker than the other, and the folds consistently chose to
down-weight the weaker member. A fusion of two comparably strong and genuinely
complementary predictors might well behave differently. Nothing here tests that.

It is also not a claim that the earlier result was dishonestly produced. The
protocol that generated it is easy to arrive at by accident, because a weight
sweep serves two different purposes that look the same on the page. Presented as
a sensitivity analysis, it answers "how fragile is this result to the weight?"
Used to pick the reported row, it answers "what is the best this combination can
look like?" The second question's answer is then reported as though it were the
first's. The diagnostic is procedural rather than statistical: ask whether the
configuration being reported was selected using the labels it is scored against.

What the result does establish is that in this corner of the literature the
difference between the two protocols is large enough to manufacture a positive
finding out of a null one, on data where the optimistic protocol produced a
plausible-looking effect with a bootstrap interval attached. Reporting an
interval does not protect against selection bias when the selection happened
before the interval was computed. Given how inexpensive the nested alternative
is --- one extra loop over
sequences, no new data, no new model --- we think it should be the default
whenever a combination parameter is tuned on a dataset of this size.

\subsection{Relation to prior work}

Our benchmark is built on the BP3C50ID split distributed with BepiPred-3.0
\citep{clifford2022bepipred}. The local predictor descends from the
feature-based and deep approaches surveyed by \citet{cia2023critical} and
exemplified in this journal by \citet{attique2023deepbce}. We do not advance the
state of the art in predictive accuracy and do not claim to. Standalone
CARBON-BCE-R remains below BepiPred-3.0 on BP3C50ID, and structure-aware methods
such as DiscoTope-3.0 \citep{hoie2024discotope} occupy ground our sequence-only
setting does not reach.

Our evaluation is small beside purpose-built benchmarks for this area such as
AsEP \citep{liu2024asep}, and we call it a measurement harness rather than a
benchmark for that reason. Within the scope of inference routing, however, we
are not aware of prior work that routes between a local and a remote
conformational epitope predictor and measures the consequences. Federated
approaches to protein modelling address a different problem. They distribute
training rather than place inference.

The systems contribution sits alongside the split-computing literature
\citep{matsubara2022split}, from which the routing mechanism is borrowed rather
than invented. What we add is the measurement of that mechanism against the
objectives a wet-lab screening workflow actually has. That includes the
counting measure of raw-sequence disclosure, which accuracy-oriented
formulations do not express, and the observation that the energy argument
inverts at this scale.

The methodological contribution is an instance of the problem
\citet{kapoor2023leakage} document across applied machine-learning fields,
reported from the inside: we are correcting our own earlier result rather than
auditing someone else's. We think instances of this kind are useful precisely
because the published record contains few of them.
